\documentclass[10pt,a4paper]{amsart}
\usepackage[margin=19mm]{geometry}
\usepackage{amsmath,amssymb,mathtools}
\usepackage[scr=rsfs,cal=euler]{mathalfa}
\usepackage{xfrac}
\usepackage[unicode,hidelinks,bookmarks=false]{hyperref}
\newcommand{\ie}{i.e.\ }
\newcommand{\cf}{cf.\ }
\newcommand{\resp}{respectively\ }
\newcommand{\Subsection}{\subsection}
\providecommand{\nobreakdash}{\nobreak\hbox{-}}
\setlength{\emergencystretch}{2em}
\multlinegap0pt
\usepackage{bm}
\usepackage{bbm}
\usepackage{dsfont}
\usepackage{xspace}
\usepackage{cancel}
\usepackage{mathtools}
\usepackage{amsthm}
\makeatletter
\@ifundefined{equation}{\newtheorem{equation}{Equation}}{}
\@ifundefined{proposition}{\newtheorem{proposition}[equation]{Proposition}}{}
\@ifundefined{theorem}{\newtheorem{theorem}[equation]{Theorem}}{}
\makeatother
\title[Linear representations, crystallographic quotients, and twis]{Linear representations, crystallographic quotients, and twisted conjugacy of virtual Artin groups}
\author{Neeraj Kumar Dhanwani, Pravin Kumar, Tushar Kanta Naik,, Mahender Singh}
\date{Source: arXiv:2409.10270v3 (CC BY 4.0)}
\begin{document}
\maketitle

\noindent\emph{Source and licence.} Linear representations, crystallographic quotients, and twisted conjugacy of virtual Artin groups. Version arXiv:2409.10270v3 (CC BY 4.0), distributed under the Creative Commons Attribution 4.0 International licence, as recorded in the arXiv licence metadata for this exact version and shown on the abstract page https://arxiv.org/abs/2409.10270v3. Full rights notice: licenses/arxiv-2409.10270-CC-BY-4.0.txt.\par\smallskip

\noindent\textit{Editorial scope.} Counted unit: the Theorem whose source label is \texttt{crystallographic conjugacy} in the article (source0.tex lines 592--596, source label \texttt{crystallographic conjugacy}), delivered together with the article's own complete original proof of it (source0.tex lines 597--710, one \texttt{proof} environment of 114 lines). No mathematics is written, repaired or re-derived: statement and proof are the article's own text line for line. Required context retained uncounted: prop:freeabelian (lines 494--496). Every definition, display and label of those lines is retained unchanged. Omitted from this package and not redistributed: 1--493, 497--591, 711--826. Nothing inside a retained excerpt is omitted or altered: every retained body is delivered exactly as the article's lines are written, so no omission receipt of any kind accompanies this package. Citations: the retained excerpts carry no citation command at all, so no bibliography entry of the article is transcribed and no reference is redistributed. Reference adaptation: none was needed and none was made; every reference inside the delivered unit resolves inside it, so no unresolved reference or citation is produced. Printed slips kept literal: no printed word, symbol, hypothesis or number of the counted statement or of its proof was altered. Layout and class: the article's own class is \texttt{amsart} with packages including amscd, amsfonts, amsmath, amssymb, amsthm, mathrsfs, appendix, enumerate, graphicx, flexisym, epsfig, float, hyperref, pdfpages; the wrapper is the lane's pinned \texttt{amsart} editorial wrapper at 10pt on A4, which restates the theorem-like environments the retained text needs and adds the article's own macro definitions that the retained text needs (none), with extra wrapper packages (bm, bbm, dsfont, xspace, cancel, mathtools). The delivered numbering is the unit's own. Assets: no retained span contains a figure environment, a graphics-inclusion command, a PDF-page inclusion, a table or a TikZ picture; no image file is copied into this package, the package declares no asset dependency and no image file is redistributed with it. Licence: version arXiv:2409.10270v3 of this article is distributed under the Creative Commons Attribution 4.0 International licence, as recorded in the arXiv licence metadata for this exact version and shown on the abstract page https://arxiv.org/abs/2409.10270v3. A full rights notice is delivered at licenses/arxiv-2409.10270-CC-BY-4.0.txt. Characters: every retained excerpt is pure ASCII, so the delivered engine, XeLaTeX with its default Latin Modern text font, reports no missing character.
\begin{proposition}\label{prop:freeabelian}
Let $VA[\Gamma]$ be the virtual Artin group associated with a Coxeter graph $\Gamma$, and let $\Phi[\Gamma]$ be it root system. Then the group $PVA[\Gamma]/PVA[\Gamma]'$ is a free abelian group of rank $|\Phi[\Gamma]|$.
\end{proposition}

\begin{theorem}\label{crystallographic conjugacy}
Let $VA[\Gamma]$ be a spherical virtual Artin group. The elements $w=(\prod_{\beta \in \Phi[\Gamma]} \overline{\zeta}_{\beta}^{\,a_\beta})\theta$ and $w'=(\prod_{\beta \in \Phi[\Gamma]} \overline{\zeta}_{\beta}^{ \,b_\beta})\theta$ are conjugate in $VA[\Gamma]/PVA[\Gamma]'$ if and only if 
$$\sum_{0\le t \le m_\beta-1}f_{\rho(\theta)^{t}(\beta)}+c_{\rho(\overline{\pi}_P(\eta))^{-1}(\beta)}=d_\beta$$
where $\zeta_\beta \in T_\theta$, $m_\beta=|\mathcal{O}_{\theta}(\zeta_{\beta})|$, $c_{\gamma}=\sum_{\zeta_\beta\in\mathcal{O}_{\theta}(\zeta_{\gamma})}{a_\beta}$, $d_{\gamma}=\sum_{\zeta_\beta\in\mathcal{O}_{\theta}(\zeta_{\gamma})}{b_\beta}$, $c_{\rho(\overline{\pi}_P(\eta))^{-1}(\gamma)}=\sum_{\zeta_\beta\in\mathcal{O}_{\theta}(\zeta_{\rho(\overline{\pi}_P(\eta))^{-1}(\gamma)})}{a_\beta}$ for $\zeta_\gamma \in T_\theta$ and the element  $\eta\in VA[\Gamma]/PVA[\Gamma]'$ is chosen such that $\overline{\pi}_P(\eta)$ is in the centralizer of $\theta$ in $W[\Gamma]$ with $\eta\theta\eta^{-1}= (\prod_{\beta\in \Phi[\Gamma]} \overline{\zeta}_{\beta}^{\,f_{\beta}})\theta$.
\end{theorem}

\begin{proof}
Let $T_\theta=\{\zeta_{\beta_1}, \ldots, \zeta_{\beta_r}\}$ be a set of representatives of orbits of the action of $\theta$ on $\{\zeta_\beta ~|~ \beta \in \Phi[\Gamma]\}$. We claim that $w$ is conjugate to  $(\prod_{\zeta_{\beta_i}\in T_{\theta}} \overline{\zeta}_{\beta_i}^{\,c_{\beta_i}})\theta$, where $c_{\beta_i}=\sum_{\zeta_\beta\in\mathcal{O}_{\theta}(\zeta_{\beta_i})}{a_\beta}$ for each $\zeta_{\beta_i}\in T_{\theta}$, whereas $w'$ is conjugate to $(\prod_{\zeta_{\beta_i}\in T_{\theta}} \overline{\zeta}_{\beta_i}^{\,d_{\beta_i}})\theta$, where $d_{\beta_i}=\sum_{\zeta_\beta\in\mathcal{O}_{\theta}(\zeta_{\beta_i})}{b_\beta}$ for each $\zeta_{\beta_i}\in T_{\theta}$
\par

Let $Y=\prod_{\beta \in \Phi[\Gamma]}\overline{\zeta}_{\beta}^{\, y_\beta}\in PVA[\Gamma]/PVA[\Gamma]'$, where $y_\beta\in\mathbb Z$. Then we have
\begin{eqnarray*}
& &Y \Bigg(\big(\prod_{\beta \in \Phi[\Gamma]}\overline{\zeta}_{\beta}^{\,a_\beta} \big)\theta \Bigg)Y^{-1}=\big(\prod_{\zeta_{\beta_i}\in T_{\theta}} \overline{\zeta}_{\beta_i}^{\,c_{\beta_i}} \big)\theta\\
\Leftrightarrow & & Y \Bigg( \big(\prod_{\beta \in \Phi[\Gamma]} \overline{\zeta}_{\beta}^{\,a_\beta} \big)\theta  \Bigg)Y^{-1} \theta^{-1}=\prod_{\zeta_{\beta_i}\in T_{\theta}} \overline{\zeta}_{\beta_i}^{\,c_{\beta_i}}\\
\Leftrightarrow & & Y \Bigg(\prod_{\beta \in \Phi[\Gamma]} \overline{\zeta}_{\beta}^{\,a_\beta}\Bigg)  \Bigg(\prod_{\beta \in \Phi[\Gamma]}\overline{\zeta}_{{\rho(\theta)(\beta)}}^{\,-y_\beta} \Bigg)=\prod_{\zeta_{\beta_i}\in T_{\theta}}\overline{\zeta}_{\beta_i}^{\,c_{\beta_i}}
\end{eqnarray*}
\begin{align}\label{se1}
\Leftrightarrow\begin{cases}
	y_{\beta}+a_{\beta}- y_{\rho(\theta)^{-1}(\beta)}=0 &~\mbox{ if }~\zeta_\beta\notin T_{\theta},\\
	y_{\beta_i}+a_{\beta_i}-y_{\rho(\theta)^{-1}(\beta_i)}=c_{\beta_i} &~\mbox{ if }~\zeta_\beta=\zeta_{\beta_i}\in T_{\theta}.
\end{cases}
\end{align}
The system of equations \eqref{se1} has a solution if and only if the following subsystem of equations has a solution   
	\begin{align}\label{se2}
\begin{cases}
		y_{\rho(\theta)^t(\beta_i)}+a_{\rho(\theta)^t(\beta_i)}-y_{\rho(\theta)^{-(1+m_{i}-t)}(\beta_i)}=0,\\
		y_{\beta_i}+a_{\beta_i}-y_{\rho(\theta)^{-1}(\beta_i)}=c_{\beta_i},
	\end{cases}
\end{align}
for all $1\le t<m_{i}$ and $\zeta_{\beta_i}\in T_{\theta}$, where $m_{i}=|\mathcal{O}_{\theta}(\zeta_{\beta_i})|$.  We show that the subsystems of equations \eqref{se2} admits a solution. For fixed $\zeta_{\beta_i}\in T_{\theta}$, consider the elements $e_1=\zeta_{\beta_i},e_2=\zeta_{\rho(\theta)(\beta_i)},\dots,e_{m_{i}}=\zeta_{\rho(\theta)^{m_{i}-1}(\beta_i)}$. Then the  matrix of $\theta$ with respect to the ordered set $\{e_1,e_2,\dots,e_{m_i}\}$ is given by
$$M_{\beta_i}=
\left(\begin{matrix}	0&0&\ldots&0&1\\
	1&0&\ldots&0&0\\
	0&1&\ldots&0&0\\
	\vdots&\vdots&\ddots&\vdots&\vdots\\ 
	0&0&\ldots&1&0
\end{matrix}\right).
$$
Thus, for each $\zeta_{\beta_i}\in T_{\theta}$, the subsystems of equations \eqref{se2} can be written as
$$
[I_{\beta_i}-M_{\beta_i}]
\left(\begin{matrix}
	y_{\beta_i}\\
	y_{\rho(\theta)(\beta_i)}\\
	\vdots\\
	y_{\rho(\theta)^{m_{i}-1}(\beta_i)}
\end{matrix}\right)=
\left(\begin{matrix}	-a_{\beta_i}+c_{\beta_i}\\
	-a_{\rho(\theta)(\beta_i)}\\
	\vdots\\
	-a_{\rho(\theta)^{m_{i}-1}(\beta_i)}
\end{matrix}\right),
$$ 
where $I_{\beta_i}$ is the $m_{i}\times m_{i}$ identity matrix and $m_i=|\mathcal{O}_{\theta}(\zeta_{\beta_i})|$. Analysing the matrix, we conclude that \eqref{se2}) has a solution if and only if $$c_{\beta_i}=\sum_{\zeta_\beta\in\mathcal{O}_{\theta}(\zeta_{\beta_i})}{a_\beta}=\sum_{0 \le \ell \le m_{i}-1} a_{\rho(\theta)^\ell(\beta_i)}$$ for all $\zeta_{\beta_i}\in T_{\theta}$. Similar arguments establish the second assertion of the claim. 
\par

Next, we analyse the conditions under which  $(\prod_{\zeta_{\beta_i}\in T_{\theta}} \overline{\zeta}_{\beta_i}^{\, c_{\beta_i}})\theta$ and $(\prod_{\zeta_{\beta_i}\in T_{\theta}} \overline{\zeta}_{\beta_i}^{\, d_{\beta_i}})\theta$ are conjugate in $VA[\Gamma]/PVA[\Gamma]'$. We choose $\eta\in VA[\Gamma]/PVA[\Gamma]'$ such that $\overline{\pi}_P(\eta)$ is in the centralizer of $\theta$ in $W[\Gamma]$. This gives $\eta\theta\eta^{-1}= (\prod_{\beta\in \Phi[\Gamma]} \overline{\zeta}_{\beta}^{\, f_{\beta}})\theta$ for some $f_{\beta_i} \in \mathbb{Z}$. Let $Y=\prod_{\beta \in \Phi[\Gamma]}\overline{\zeta}_{\beta}^{\,y_\beta}\in PVA[\Gamma]/PVA[\Gamma]'$, where $y_\beta\in\mathbb Z$. Then we have
\begin{small}
\begin{eqnarray}
\nonumber & & Y\eta \Bigg(\prod_{\zeta_{\beta_i}\in T_{\theta}} \overline{\zeta}_{\beta_i}^{\, c_{\beta_i}} \Bigg) \theta\eta^{-1}Y^{-1}=(\prod_{\zeta_{\beta_i}\in T_{\theta}} \overline{\zeta}_{\beta_i}^{\,d_{\beta_i}})\theta\\
\nonumber  \Leftrightarrow & & Y\eta \Bigg(\prod_{\zeta_{\beta_i}\in T_{\theta}}\overline{\zeta}_{\beta_i}^{\,c_{\beta_i}}\Bigg) \theta\eta^{-1}Y^{-1}\theta^{-1}=\prod_{\zeta_{\beta_i}\in T_{\theta}} \overline{\zeta}_{\beta_i}^{\, d_{\beta_i}}\\
\nonumber  \Leftrightarrow & & Y \Bigg(\prod_{\zeta_{\beta_i}\in T_{\theta}} \overline{\zeta}_{{\rho(\overline{\pi}_P(\eta))(\beta_i)}}^{\,c_{\beta_i}}\Bigg) \eta\theta\eta^{-1}Y^{-1}\theta^{-1}=\prod_{\zeta_{\beta_i}\in T_{\theta}} \overline{\zeta}_{\beta_i}^{\,d_{\beta_i}} ,\\
\nonumber  \Leftrightarrow & & Y \Bigg(\prod_{\zeta_{\beta_i}\in T_{\theta}} \overline{\zeta}_{{\rho(\overline{\pi}_P(\eta))(\beta_i)}}^{\, c_{\beta_i}}\Bigg) \Bigg(\prod_{\beta \in \Phi[\Gamma]} \overline{\zeta}_{\beta}^{\,f_\beta} \Bigg) \theta Y^{-1}\theta^{-1}=\prod_{\zeta_{\beta_i}\in T_{\theta}} \overline{\zeta}_{\beta_i}^{\,d_{\beta_i}},\\
\label{mid equation} \quad \Leftrightarrow & & \Bigg(\prod_{\beta \in \Phi[\Gamma]}\overline{\zeta}_{\beta}^{\,y_\beta} \Bigg) \Bigg(\prod_{\zeta_{\beta_i} \in T_{\theta}} \overline{\zeta}_{{\rho(\overline{\pi}_P(\eta))(\beta_i)}}^{\,c_{\beta_i}} \Bigg) \Bigg(\prod_{\beta \in \Phi[\Gamma]} \overline{\zeta}_{\beta}^{\,f_\beta} \Bigg) \Bigg(\prod_{\beta \in \Phi[\Gamma]}\overline{\zeta}_{{\rho(\theta)(\beta)}}^{\, -y_\beta} \Bigg)=\prod_{\zeta_{\beta_i}\in T_{\theta}} \overline{\zeta}_{\beta_i}^{\,d_{\beta_i}}.
\end{eqnarray}
\end{small}


We have defined $c_{\beta_i}$ for each $\zeta_{\beta_i} \in T_{\theta}$. Given any $\beta \in \Phi[\Gamma]$, we have $\zeta_\beta\in  \mathcal{O}_{\theta}(\zeta_{\beta_i})$ for some $\zeta_{\beta_i} \in T_\theta$. In this case, we set $c_\beta=c_{\beta_{i}}$.  For each $\zeta_{\beta_i}\in T_\theta$, note that $\rho(\overline{\pi}_P(\eta))^{-1}\cdot\zeta_{\beta_i}=\zeta_{\rho(\overline{\pi}_P(\eta))^{-1}(\beta_i)} \in \mathcal{O}_{\theta}(\zeta_{\beta_j})$ for a unique $\zeta_{\beta_j} \in T_\theta$. This implies that there exists an integer $0\leq l'\leq m_j$ such that  $\rho(\theta)^{l'}\cdot \zeta_{\beta_j}=\rho(\overline{\pi}_P(\eta))^{-1}\cdot\zeta_{\beta_i}$. Since  $\overline{\pi}_P(\eta)$ is in the centralizer of $\theta$, we conclude that $\zeta_{\rho(\overline{\pi}_P(\eta))(\beta_j)}=\zeta_{\rho(\theta)^{-l'}(\beta_i)}$. Setting $l = -l' \mod m_i$, we obtain $\zeta_{\rho(\overline{\pi}_P(\eta))(\beta_j)}=\zeta_{\rho(\theta)^{l}(\beta_i)}$. In view of Proposition \ref{prop:freeabelian}, this further gives  $\rho(\overline{\pi}_P(\eta))(\beta_j)=\rho(\theta)^l(\beta_i)$.
\par

It is easy to see that $\zeta_{\beta_j}\in \mathcal{O}_{\theta}(\zeta_{\rho(\overline{\pi}_P(\eta))^{-1}(\beta_i)})$, and hence $c_{\beta_j} = c_{\rho(\overline{\pi}_P(\eta))^{-1}(\beta_i)}$. If $l=0$, then equation \eqref{mid equation} holds if and only if the following system of equations has a solution 
\begin{align}\label{mid equation 1}
\begin{cases}
y_{\rho(\theta)^k(\beta_i)}+f_{\rho(\theta)^k(\beta_i)}-y_{\rho(\theta)^{-(1+m_i-k)}(\beta_i)}=0,\\
y_{\beta_i}+f_{\beta_i}+c_{\rho(\overline{\pi}_P(\eta))^{-1}(\beta_i)}-y_{\rho(\theta)^{-1}(\beta_i)}=d_{\beta_i},	
\end{cases} 
\end{align}
for $1\le k<m_i$, where $m_i=|\mathcal{O}_{\theta}(\zeta_{\beta_i})|$. If $1\leq l \leq m_i$, then equation \eqref{mid equation} holds if and only if the following system of equations has a solution 
\begin{align}\label{mid equation 2}
\begin{cases}
y_{\rho(\theta)^k(\beta_i)}+f_{\rho(\theta)^k(\beta_i)}-y_{\rho(\theta)^{-(1+m_i-k)}(\beta_i)}=0,\\
y_{\rho(\theta)^l(\beta_i)}+f_{\rho(\theta)^l(\beta_i)}+c_{\rho(\overline{\pi}_P(\eta))^{-1}(\beta_i)}-y_{\rho(\theta)^{-(1+m_i-l)}(\beta_i)}=0,\\
y_{\beta_i}+f_{\beta_i}-y_{\rho(\theta)^{-1}(\beta_i)}=d_{\beta_i},	
\end{cases} 
\end{align}
for $1\le k<m_i$ and $k\not=l$, where $m_i=|\mathcal{O}_{\theta}(\zeta_{\beta_i})|$.  For fixed $\zeta_{\beta_i}\in T_{\theta}$, if $l=0$, then the system of equations \eqref{mid equation 1} can be written as
$$
[I_{\beta_i}-M_{\beta_i}]
\left(\begin{matrix}
	y_{\beta_i}\\
	y_{\rho(\theta)(\beta_i)}\\
	\vdots\\
	y_{\rho(\theta)^{m_{i}-1}(\beta_i)}
\end{matrix}\right)=
\left(\begin{matrix}	-f_{\beta_i}-c_{\rho(\overline{\pi}_P(\eta))^{-1}(\beta_i)}+d_{\beta_i}\\
	-f_{\rho(\theta)(\beta_i)}\\
	\vdots\\
	-f_{\rho(\theta)^{m_{i}-1}(\beta_i)}
\end{matrix}\right).
$$ 
Similarly, if $1\leq l \leq m_i$, then the system of equations \eqref{mid equation 2} can be written as
$$
[I_{\beta_i}-M_{\beta_i}]
\left(\begin{matrix}
	y_{\beta_i}\\
	\vdots\\
		y_{\rho(\theta)^l(\beta_i)}\\
	\vdots\\
	y_{\rho(\theta)^{m_{i}-1}(\beta_i)}
\end{matrix}\right)=
\left(\begin{matrix}	-f_{\beta_i}+d_{\beta_i}\\
	\vdots\\
		-f_{\rho(\theta)^l(\beta_i)}-c_{\rho(\overline{\pi}_P(\eta))^{-1}(\beta_i)}\\
	\vdots\\
	-f_{\rho(\theta)^{m_{i}-1}(\beta_i)}
\end{matrix}\right).
$$ 
Analysing the matrices, we conclude that \eqref{mid equation 1} and \eqref{mid equation 2} have a solution if and only if $\sum_{0\le t \le m_i-1}f_{\rho(\theta)^{t}(\beta_i)}+c_{\rho(\overline{\pi}_P(\eta))^{-1}(\beta_i)}=d_{\beta_i}$ for all $\zeta_{\beta_i}\in T_{\theta}$. This completes the proof of the theorem.
\end{proof}

\end{document}
